%% ============================================================
%% cap07.tex — Capítulo 7
%% Software generativo, ética y aprendizaje de la programación
%% Autor: Roberto Albeiro Pava-Díaz, Ph.D.
%% Universidad Distrital Francisco José de Caldas
%% Versión Phase 2 — Marzo 2026
%% ============================================================

\chapter{Software generativo, ética y aprendizaje de la programación}
\label{cap:ia-etica}


\noindent
Once \textit{sprints} después del primero, el estudiante llega a este
capítulo con código escrito, diagramas producidos, pruebas ejecutadas
y, muy probablemente, decenas de interacciones con asistentes
computacionales generativos a lo largo del proceso. La pregunta que
se planteó en la advertencia del prefacio ---si el empleo de software
generativo fortalece o menoscaba el aprendizaje genuino de la
\poo--- ya no es abstracta: hay evidencia propia con la que responderla.
Queda, sin embargo, la duda de si esa evidencia es suficiente para
generalizar, y el autor reconoce que no existe, al momento de redactar
este texto (marzo de 2026), un cuerpo empírico robusto que resuelva
la pregunta para el contexto específico de los estudiantes de pregrado
de ingeniería en universidades públicas colombianas.\footnote{%
  La revisión de la literatura realizada en enero de 2026 en las bases
  de datos Scopus, IEEE Xplore y ACM Digital Library arrojó escasos
  estudios controlados sobre el impacto del software generativo en el
  aprendizaje de la POO en contextos latinoamericanos. La mayoría de los
  trabajos disponibles provienen de universidades norteamericanas o europeas,
  con estudiantes cuyas condiciones de acceso tecnológico difieren
  sustancialmente de las prevalentes en la \ud. La extrapolación exige cautela.}

Tres preguntas organizan la arquitectura de este capítulo final.
Primera: ¿qué son, técnicamente, estos asistentes computacionales
generativos y qué pueden o no pueden hacer en el dominio de la
programación orientada a objetos? Segunda: ¿qué disposiciones pedagógicas
y éticas se derivan del hecho de que el estudiante tenga acceso a ellos?
Tercera: ¿cómo puede el formador evaluar el razonamiento de diseño
orientado a objetos en un entorno donde el código puede originarse en
cualquier fuente? Ninguna de estas preguntas admite respuesta definitiva.
Pero todas admiten un análisis más riguroso del que suele encontrarse en
el discurso ordinario sobre IA y educación.

%% ─────────────────────────────────────────────
\section{Modelos generativos en el contexto educativo}
\label{sec:modelos-generativos}
%% ─────────────────────────────────────────────

\noindent
Dos mil veintitrés fue el año en que los asistentes computacionales
generativos dejaron de ser una curiosidad académica para convertirse
en presencia cotidiana en los laboratorios de programación. Dos mil
veintiséis es el año en que este texto se escribe, y la penetración
de estos sistemas en los entornos de aprendizaje universitario ---no
solo en Colombia, sino en América Latina y el mundo--- se ha
profundizado de manera que pocos anticiparon con precisión.\footnote{%
  El autor observó, durante el primer semestre de 2025 en la
  Facultad de Ingeniería de la \ud, que más del 70\,\% de los
  estudiantes de Programación Orientada a Objetos declararon
  voluntariamente haber empleado algún asistente computacional
  generativo en la elaboración de al menos un entregable del semestre.
  La cifra proviene de una encuesta anónima aplicada al cierre del
  semestre y no ha sido sometida a validación estadística rigurosa;
  se reporta como indicador orientador, no como medida definitiva.}
Conviene, antes de examinar las implicaciones pedagógicas y éticas,
caracterizar con precisión qué son estos sistemas y cómo operan.

\subsection{Arquitectura conceptual: transformadores y predicción de siguiente \textit{token}}
\label{subsec:arquitectura-llm}

Los sistemas que generan texto y código de manera automática ---denominados
en este libro \emph{asistentes computacionales generativos} o, con mayor
brevedad, \emph{software generativo}--- se fundamentan en una arquitectura
neuronal denominada \textit{transformer}, propuesta por Vaswani et al.\ en 2017
y perfeccionada de manera acelerada durante la década siguiente
\parencite{Martin2017}. El principio de funcionamiento, descrito a nivel
conceptual, es la predicción iterativa del siguiente \textit{token} más probable
dada una secuencia de entrada: el modelo no ``comprende'' el código en ningún
sentido semántico profundo; estima la continuación estadísticamente más verosímil
de una cadena de símbolos. Esta distinción ---entre plausibilidad estadística y
comprensión semántica--- no es menor. Es, de hecho, la fuente de la mayoría de
los errores que estos sistemas cometen cuando se les solicita código que modela
dominios específicos con invariantes no triviales.

El ajuste por retroalimentación humana ({\small RLHF}, por sus siglas en inglés
\textit{Reinforcement Learning from Human Feedback}) permite orientar el
comportamiento del modelo hacia respuestas evaluadas como útiles por revisores
humanos \parencite{Fowler2018}. El resultado es un sistema que produce texto y código de alta fluidez
superficial pero con garantías semánticas limitadas. Queda por determinar si
el RLHF introduce o reduce sesgos específicos en el dominio de la POO; los
indicios disponibles al momento de redactar este texto sugieren que los sesgos
existen pero no han sido caracterizados con precisión para dominios institucionales
del tipo del estatuto de la \ud.

Cinco propiedades de estos modelos resultan directamente relevantes para el
diseño de estrategias pedagógicas:

\begin{enumerate}[label=\arabic*.]
  \item \textbf{Alta fluidez superficial}: el código generado compila con
        frecuencia y luce correcto a primera vista, lo que dificulta
        la detección de errores semánticos.
  \item \textbf{Dependencia del corpus de entrenamiento}: la calidad del
        código generado para un lenguaje o dominio específico depende
        del volumen y la calidad del código de ese lenguaje en el corpus
        de entrenamiento.
  \item \textbf{Incapacidad para razonar sobre invariantes no documentadas}:
        si una restricción del dominio no está explícitamente en el
        \textit{prompt}, el modelo la ignorará.
  \item \textbf{Alucinación}: el modelo puede generar código que invoca
        métodos o clases inexistentes en la API objetivo, con
        confianza aparente.
  \item \textbf{Genericidad de los patrones}: el modelo tiende a producir
        soluciones genéricas que pueden violar los principios SOLID
        cuando el dominio exige especialización.
\end{enumerate}

\subsection{Taxonomía de asistentes computacionales generativos para programación}
\label{subsec:taxonomia-asistentes}

La Tabla~\ref{tab:asistentes-generativos} clasifica los principales asistentes
computacionales generativos disponibles al primer trimestre de 2026, diferenciando
entre los de propósito general y los especializados en código. La clasificación es
necesariamente provisional: el ritmo de evolución de estos sistemas hace que
cualquier fotografía del ecosistema tenga vigencia constreñida.

\begingroup
\small
\begin{longtable}{L{2.5cm} L{2cm} L{4.5cm} L{4.5cm}}
	% --- Título y encabezado de la primera página ---
	\caption{Taxonomía de asistentes computacionales generativos relevantes para el aprendizaje de la POO. Elaboración del autor (primer trimestre de 2026).} \label{tab:asistentes-generativos} \\
	\toprule
	\textbf{Asistente} & \textbf{Tipo} & \textbf{Capacidades destacadas} & \textbf{Limitaciones documentadas} \\
	\midrule
	\endfirsthead
	
	% --- Encabezado para las páginas siguientes ---
	\multicolumn{4}{c}{{\bfseries \tablename\ \thetable{} -- Continuación}} \\
	\toprule
	\textbf{Asistente} & \textbf{Tipo} & \textbf{Capacidades destacadas} & \textbf{Limitaciones documentadas} \\
	\midrule
	\endhead
	
	% --- Pie de tabla para páginas intermedias ---
	\midrule
	\multicolumn{4}{r}{{Continúa en la siguiente página...}} \\
	\bottomrule
	\endfoot
	
	% --- Pie de tabla final ---
	\bottomrule
	\endlastfoot
	
	% --- Datos de la tabla ---
	GitHub Copilot
	& Especializado en código
	& Autocompletado en contexto del IDE; sugerencias de métodos completos; integración con repositorios propios del usuario
	& Dependencia del contexto visible en el editor; errores en código que modela dominios complejos con múltiples invariantes; código \java\ frecuentemente genérico \\
	\midrule
	
	Claude (Anthropic)
	& Propósito general (código fuerte)
	& Razonamiento sobre requisitos complejos; explicación detallada de decisiones de diseño; análisis crítico del código aportado
	& Alucinación de API en versiones de JDK recientes; restricciones en ventana de contexto para proyectos grandes \\
	\midrule
	
	ChatGPT (OpenAI)
	& Propósito general
	& Generación de código en múltiples lenguajes; explicación de mensajes de error; generación de pruebas unitarias
	& Patrones de diseño genéricos; dificultad con dominios institucionales específicos no documentados en el corpus \\
	\midrule
	
	Gemini (Google)
	& Propósito general / Multimodal
	& Integración con Google Workspace; razonamiento multimodal; explicación de diagramas UML en imagen
	& Menor especialización en código \java\ empresarial comparado con Copilot; variabilidad en sintaxis de \rust \\
	\midrule
	
	Cursor / Aider
	& Especializados en IDE
	& Edición de archivos completos en contexto; refactorización guiada; \textit{diff} aplicado directamente al código base
	& Requieren configuración avanzada; curva de aprendizaje para estudiantes sin experiencia previa en CLI \\
	\midrule
	
	Codeium / Tabnine
	& Especializados (opciones locales)
	& Autocompletado gratuito; menor latencia; pueden ejecutarse con modelos locales en entornos con baja conectividad
	& Capacidad de razonamiento menor que modelos de mayor escala; ventana de contexto limitada \\
\end{longtable}
\endgroup

\subsection{Penetración en los entornos universitarios colombianos y latinoamericanos}
\label{subsec:penetracion-col}

La difusión del software generativo en los programas de ingeniería de sistemas
de Colombia presenta una asimetría que el autor denomina ---quizás de manera
provisional--- \emph{brecha generativa}: la diferencia entre las condiciones
de acceso a estos sistemas entre instituciones privadas de alta matrícula y las
universidades públicas. La \ud, con sede en la Macarena (carrera 8 con calle 40)
y salas de sistemas compartidas cuyo ancho de banda es insuficiente en horas
pico, no puede asumir que todos sus estudiantes acceden a versiones premium de
estos asistentes desde dispositivos propios con conexión estable. Esa
constatación ---que no es un detalle menor, sino una condición estructural del
aprendizaje--- atraviesa toda la discusión ética de la sección~\ref{sec:marco-etico}.

Sin embargo, la penetración existe. Cabría pensar que la ausencia de acceso
premium implica ausencia de uso; la experiencia del autor en el primer semestre
de 2025 contradice esa hipótesis. Los estudiantes emplean versiones gratuitas,
acuden a compañeros con acceso de pago, utilizan los quince minutos diarios
de uso gratuito con múltiples cuentas. La creatividad para sortear las barreras
de acceso es, en sí misma, un dato pedagógicamente relevante: el estudiante
distrital no espera que la herramienta llegue a él; la busca. No es descartable
que esa búsqueda activa contribuya a un uso más reflexivo que el de quien tiene
acceso ilimitado sin esfuerzo previo.

A nivel latinoamericano, el panorama presenta similitudes. Los estudios
preliminares disponibles ---entre ellos, encuestas de adopción en universidades
brasileñas, mexicanas y argentinas publicadas entre 2024 y 2025--- indican
tasas de adopción declarada superiores al 60\,\% entre estudiantes de ingeniería
de sistemas, con variaciones significativas según el nivel del curso y la
actitud del docente \parencite{Bloch2018}. La mayoría de las instituciones
latinoamericanas carece, al momento de redactar este texto, de políticas
explícitas sobre el uso de software generativo en evaluaciones de programación.

%% ─────────────────────────────────────────────
\section{Oportunidades y riesgos del uso de asistentes computacionales en \poo}
\label{sec:oportunidades-riesgos}
%% ─────────────────────────────────────────────

\noindent
El análisis equilibrado del software generativo en el contexto del aprendizaje
de la orientación a objetos exige resistir dos tentaciones simétricas: la del
entusiasmo acrítico, que ve en estos sistemas una revolución pedagógica sin
precedentes, y la del rechazo apriorístico, que los descarta como amenaza para
la formación genuina. Ninguna de las dos posiciones resiste el escrutinio de
la evidencia disponible. Lo que se propone aquí es un examen caso a caso,
anclado en el dominio específico del sistema de información estatutario
que atraviesa el libro.

\subsection{Oportunidades: andamiaje, explicación y comparación}
\label{subsec:oportunidades}

Cuatro oportunidades concretas se identifican, cada una con su ámbito de validez
y sus condiciones de uso:

\textbf{Primera: generación de código de andamiaje} (\textit{boilerplate}).
En los Sprints 1 a 3 (capítulos~\ref{cap:fundamentos} y~\ref{cap:modelado}),
el estudiante dedica tiempo considerable a escribir constructores, métodos
\textit{getters} y \textit{setters}, y métodos \texttt{toString} y \texttt{equals}
en \java, o las implementaciones de \textit{traits} \texttt{Display} y \texttt{Debug}
en \rust. El software generativo puede producir este código estructural de
manera confiable ---puesto que se trata de patrones altamente repetidos en el
corpus de entrenamiento--- liberando tiempo cognitivo del estudiante para las
decisiones de diseño que sí requieren comprensión del dominio. Esta es, quizás,
la oportunidad de uso más legítima y menos arriesgada.

\textbf{Segunda: explicación de mensajes de error del compilador.}
Los mensajes del compilador de \rust, en particular, son notoriamente
detallados pero inicialmente opacos para el estudiante novel. Preguntar
al asistente computacional «¿qué significa este error de \textit{lifetime}?»
y obtener una explicación contextualizada produce un valor pedagógico genuino
---siempre que el estudiante entienda la explicación antes de aplicarla, y no
simplemente copie la corrección sugerida sin comprenderla. La diferencia entre
estos dos modos de uso es cardinal y, con frecuencia, invisible desde afuera.

\textbf{Tercera: generación de variantes para comparación crítica.}
Una práctica que el autor ha experimentado con resultados prometedores consiste
en solicitar al asistente generativo dos o tres implementaciones alternativas
de un mismo método ---por ejemplo, tres formas distintas de representar la jerarquía
de unidades académicas del Estatuto--- y pedir al estudiante que compare
críticamente cuál es más coherente con los principios SOLID y por qué. El
software generativo pasa así de ser el productor del código a ser el generador
del problema que el estudiante debe resolver.

\textbf{Cuarta: andamiaje estructural para la reflexión.}
En la actividad final del Sprint~12 (sección~\ref{act:ia-sprint12}), el asistente
puede ayudar al estudiante a estructurar su reflexión crítica: sugerir un índice,
señalar incoherencias argumentativas, proponer preguntas que no se han formulado.
Pero el contenido ---la evidencia propia, el análisis, la posición ética--- debe
provenir del estudiante.

\subsection{Riesgos: delegación del diseño, aceptación acrítica y pérdida del proceso de depuración}
\label{subsec:riesgos}

Los riesgos, análogamente, son concretos y no abstractos:

\textbf{Primero: delegación de las decisiones de diseño orientado a objetos.}
En el Sprint~4 (capítulo~\ref{cap:java}), el estudiante enfrenta la decisión
de modelar la relación entre \texttt{Docente} y \texttt{UnidadAcademica}: ¿composición
o asociación? ¿interfaz o clase abstracta? Si delega esa decisión al asistente
sin comprenderla, el resultado puede ser código que compila pero que no refleja
las restricciones del Estatuto \parencite{Martin2017}. El riesgo no es el código malo: es el razonamiento
ausente, que permanecerá ausente también en los proyectos profesionales subsecuentes.

\textbf{Segundo: aceptación acrítica de código que viola invariantes del dominio.}
El asistente no conoce el \acuerdo\ a menos que se le proporcione. En ausencia
de ese contexto, generará código genérico que puede violar restricciones normativas
específicas. La Sección~\ref{sec:analisis-critico} documenta tres ejemplos
concretos de este riesgo. No es trivial.

\textbf{Tercero: menoscabo del proceso de depuración como fuente de aprendizaje.}
Leer el \textit{stack trace}, entender el mensaje de excepción, navegar el código
para localizar la causa raíz de un fallo: ese proceso es pedagógicamente valioso
porque obliga al estudiante a construir un modelo mental del flujo de ejecución
del programa. Si el asistente resuelve el error antes de que el estudiante lo
haya comprendido, ese aprendizaje no ocurre. La velocidad del asistente es,
en este contexto, un obstáculo.

La Tabla~\ref{tab:oportunidades-riesgos} sintetiza el análisis con ejemplos
extraídos del sistema de información estatutario.

\begingroup
\small
\begin{longtable}{L{3.5cm} L{4.5cm} L{4cm} C{1.8cm}}
	% --- Título y encabezado de la primera página ---
	\caption{Oportunidades y riesgos del empleo de software generativo en el aprendizaje de la POO (Sistema Estatutario \ud). Elaboración del autor.} \label{tab:oportunidades-riesgos} \\
	\toprule
	\textbf{Oportunidad} & \textbf{Ejemplo en el sistema} & \textbf{Riesgo} & \textbf{Sprint} \\
	\midrule
	\endfirsthead
	
	% --- Encabezado para las páginas siguientes ---
	\multicolumn{4}{c}{{\bfseries \tablename\ \thetable{} -- Continuación}} \\
	\toprule
	\textbf{Oportunidad} & \textbf{Ejemplo en el sistema} & \textbf{Riesgo} & \textbf{Sprint} \\
	\midrule
	\endhead
	
	% --- Pie de tabla para páginas intermedias ---
	\midrule
	\multicolumn{4}{r}{{Continúa en la siguiente página...}} \\
	\bottomrule
	\endfoot
	
	% --- Pie de tabla final ---
	\bottomrule
	\endlastfoot
	
	% --- Datos de la tabla ---
	Generación de \textit{boilerplate} & 
	Constructores, getters/setters y \texttt{equals()} de \texttt{Docente}, \texttt{Estudiante} y \texttt{UnidadAcademica} & 
	Delegación del diseño de herencia y jerarquías sin comprensión real & 
	2--5 \\
	\midrule
	
	Explicación de errores & 
	Mensajes de \textit{lifetime} en \rust\ al implementar préstamos entre \texttt{Facultad} y \texttt{Proyecto} & 
	Corrección mecánica sin comprensión; el error reaparece en otros contextos & 
	5--7 \\
	\midrule
	
	Generación de variantes & 
	Implementaciones de \texttt{calcularPromedio()} con enfoques iterativos, \textit{streams} y recursivos & 
	Selección de la variante más corta por estética, no por corrección semántica & 
	8 \\
	\midrule
	
	Casos de prueba & 
	Esqueletos JUnit para las invariantes de \texttt{EvaluacionFormativa} según el Estatuto & 
	Pruebas triviales que omiten las restricciones normativas reales de la \ud & 
	9--10 \\
	\midrule
	
	Andamiaje reflexivo & 
	Estructura sugerida para la reflexión crítica final del Sprint~12 & 
	Reflexión generada por IA sin evidencia de aprendizaje propio del estudiante & 
	12 \\
	\midrule
	
	Traducción de patrones & 
	Equivalente en \rust\ del patrón \textit{Builder} empleado en \java\ para \texttt{EstatutoAcademico} & 
	Código \rust\ no sintácticamente correcto que el estudiante acepta por falta de experiencia & 
	6--7 \\
\end{longtable}
\endgroup

\subsection{Tres ejemplos de código generado que viola las reglas del dominio}
\label{subsec:ejemplos-violacion}

Los tres listados que se presentan a continuación documentan situaciones reales
observadas durante el primer semestre de 2025 en las asignaturas de programación
de la Facultad de Ingeniería. En cada caso, el asistente computacional generó
código sintácticamente correcto pero semánticamente incompatible con el dominio
del \acuerdo.

\textbf{Ejemplo~1: violación de la restricción de vinculación docente.}
El Artículo~31 del \acuerdo\ establece que un docente de planta sólo puede
estar vinculado a una Facultad como unidad académica principal. El código
generado por un asistente, al no conocer esta restricción, produce:

\begin{lstlisting}[style=javaStyle,
  caption={Código generado que viola la restricción de vinculación única del
    docente (Artículo~31 del Acuerdo 004 de 2025). El error es semántico,
    no sintáctico: el compilador no lo detecta.},
  label={lst:violacion-vinculacion}]
// Código generado por asistente computacional — INCORRECTO
// Viola la restricción de vinculación única establecida en el estatuto
public class DocentePlanta extends Docente {
    // PROBLEMA: permite múltiples Facultades como afiliación principal
    private List<Facultad> facultadesAfiliadas;

    public void agregarFacultad(Facultad f) {
        // Sin validación de unicidad — viola Art. 31 del Acuerdo 004/2025
        this.facultadesAfiliadas.add(f);
    }

    public List<Facultad> getFacultades() {
        return facultadesAfiliadas;
    }
}
\end{lstlisting}

\textbf{Ejemplo~2: jerarquía de unidades académicas invertida.}
El asistente, al recibir un \textit{prompt} que menciona \texttt{Facultad},
\texttt{Proyecto Curricular} y \texttt{Escuela} sin contexto normativo,
frecuentemente invierte la jerarquía definida en el Estatuto, haciendo que
\texttt{Escuela} sea padre de \texttt{Facultad}:

\begin{lstlisting}[style=javaStyle,
  caption={Jerarquía de unidades académicas invertida por el asistente generativo.
    La relación correcta, según el Acuerdo 004 de 2025, establece que la Facultad
    contiene Proyectos Curriculares y éstos se articulan con Escuelas, no al revés.},
  label={lst:jerarquia-invertida}]
// Código generado — JERARQUÍA INCORRECTA
// El asistente asume que Escuela > Facultad, lo cual invierte
// la estructura del Estatuto de la UD

public class Escuela {
    private List<Facultad> facultades; // INCORRECTO: la Facultad no depende de Escuela
    private String nombre;

    public Escuela(String nombre) {
        this.nombre = nombre;
        this.facultades = new ArrayList<>();
    }

    public void agregarFacultad(Facultad f) {
        this.facultades.add(f); // Violación de la jerarquía normativa
    }
}
\end{lstlisting}

\textbf{Ejemplo~3: manejo de \texttt{null} donde el Estatuto exige valor explícito.}
El sistema de evaluación del \acuerdo\ no admite calificaciones nulas: toda
evaluación debe registrar un valor numérico o una ausencia justificada (\textit{enum}).
El asistente, siguiendo patrones genéricos de \java, genera código que permite
\texttt{null}:

\begin{lstlisting}[style=javaStyle,
  caption={Uso de \texttt{null} donde el dominio estatutario exige un valor
    explícito. El patrón correcto emplea \texttt{Optional<Double>} o un tipo
    enumerado que represente la ausencia justificada de calificación.},
  label={lst:null-violacion}]
// Código generado — MANEJO DE NULL INCORRECTO PARA EL DOMINIO
public class EvaluacionFormativa {
    private String codigoEstudiante;
    private String codigoAsignatura;
    private Double calificacion; // PROBLEMA: null no es válido en este dominio

    // El asistente genera un getter simple sin validación
    public Double getCalificacion() {
        return calificacion; // Puede devolver null — viola invariante del dominio
    }

    // Versión correcta exigida por el dominio:
    // public Optional<Double> getCalificacion() { ... }
    // o: private EstadoCalificacion estado; (PENDIENTE, REGISTRADA, JUSTIFICADA)
}
\end{lstlisting}

Los tres ejemplos ilustran un patrón recurrente: el asistente computacional
generativo produce código correcto para un dominio genérico imaginario, pero
incorrecto para el dominio específico del \acuerdo. La distinción no es menor;
es, de hecho, el núcleo de la discusión sobre el valor pedagógico del diseño
orientado a objetos frente a la mera producción de código.

%% ─────────────────────────────────────────────
\section{Estrategias para el análisis crítico del código generado}
\label{sec:analisis-critico}
%% ─────────────────────────────────────────────

\noindent
Dado que el software generativo produce código con fluidez superficial
pero con errores semánticos frecuentes ---tal como ilustran los listados de la
sección anterior---, resulta indispensable que el estudiante desarrolle un
protocolo sistemático de revisión crítica. No se trata de desconfiar de todo
código generado, sino de someter cada pieza a un escrutinio que el propio
estudiante debe ser capaz de realizar.

\subsection{Protocolo de revisión crítica del código generado}
\label{subsec:protocolo-revision}

El protocolo que se propone a continuación se articula en cinco dimensiones
secuenciales. Seguirlas en orden no es arbitrario: cada dimensión presupone
la anterior, y saltarse la dimensión semántica para ir directamente a la
compilabilidad equivale a aceptar que el código hace algo, sin verificar
que ese algo sea lo correcto.

\begin{enumerate}[label=\textbf{D\arabic*.}, leftmargin=2.5cm]

  \item \textbf{Corrección semántica: ¿el código hace lo que el dominio requiere?}

        Antes de ejecutar el compilador, el estudiante debe leer el código
        generado y compararlo con el requisito de dominio que motivó la
        solicitud. Para el sistema estatutario, esto implica contrastar el
        código con el artículo del \acuerdo\ correspondiente. Las preguntas
        guía son: ¿se respetan todas las restricciones normativas declaradas?
        ¿Hay decisiones de diseño que el asistente tomó de manera implícita
        sin fundamento en el Estatuto? ¿Se modelan correctamente las
        cardinalidades de las relaciones?

  \item \textbf{Corrección sintáctica y compilabilidad.}

        Sólo después de verificar la corrección semántica se procede a
        compilar. Si el código tiene errores de compilación, se entienden
        los mensajes del compilador antes de aplicar la corrección sugerida
        por el asistente. Cada error de compilación es, potencialmente, un
        evento de aprendizaje; suprimirlo sin comprenderlo es un costo
        pedagógico.

  \item \textbf{Conformidad con los principios SOLID.}

        El código generado se evalúa contra los cinco principios documentados
        en la sección~\ref{sec:solid} del capítulo~\ref{cap:fundamentos}:
        responsabilidad única, abierto/cerrado, sustitución de Liskov,
        segregación de interfaces y dependencia de abstracciones. Con
        frecuencia, el asistente viola el principio de responsabilidad única
        al agrupar en una sola clase responsabilidades que el diseño orientado
        a objetos exige separar.

  \item \textbf{Calidad de la documentación.}

        Los comentarios y la documentación generados por el asistente son
        frecuentemente genéricos. Un comentario que diga \texttt{// Returns
        the teacher's name} no añade información sobre el significado del
        nombre en el contexto del Estatuto. La revisión de la documentación
        debe verificar que los comentarios Javadoc (o los comentarios de
        documentación en \rust) sean precisos, específicos del dominio y
        referencien, cuando corresponda, el artículo normativo aplicable.

  \item \textbf{Prueba sistemática de las invariantes del dominio.}

        Cada invariante del dominio identificado en la revisión de la
        dimensión D1 debe tener al menos una prueba unitaria que lo verifique.
        No basta con que el código compile y ejecute sin error: debe demostrarse
        que las restricciones normativas se cumplen ante casos de borde
        deliberadamente diseñados para violentarlas.

\end{enumerate}

\subsection{Diagrama del protocolo de revisión}
\label{subsec:diagrama-protocolo}

El flujo del protocolo, descrito textualmente para su representación futura
como figura, puede resumirse como sigue. El punto de entrada es la recepción
del código generado. La primera verificación es semántica: si el código no
satisface los requisitos del dominio, se reformula el \textit{prompt} y se
reinicia el proceso. Si satisface los requisitos, se verifica la compilabilidad.
Si no compila, se analizan los errores antes de corrección. Si compila, se
evalúa contra SOLID. Si viola SOLID, se refactoriza. Si es conforme, se
revisa la documentación y se genera la suite de pruebas. Si las pruebas pasan,
el código es candidato a integración. Si alguna prueba falla, se regresa a D1.
El proceso es iterativo, no lineal.

\subsection{Ejemplo completo: clase \texttt{Docente} con errores ocultos y versión corregida}
\label{subsec:ejemplo-docente}

El listado~\ref{lst:docente-generado-erroneo} muestra el código generado
por un asistente computacional para la clase \texttt{Docente} ante el
\textit{prompt} siguiente: «\textit{Genera la clase Docente en Java para
un sistema de información universitario, con tipos de vinculación (planta,
cátedra, ocasional) y métodos para gestionar su asignación académica.}»
El código es sintácticamente correcto y, en apariencia, razonable. Pero
contiene al menos cuatro errores semánticos respecto del dominio del \acuerdo.

\begin{lstlisting}[style=javaStyle,
  caption={Clase \texttt{Docente} generada por asistente computacional.
    El código compila sin errores pero contiene cuatro violaciones del dominio
    estatutario, identificadas con comentarios \texttt{// ERROR-n:}.},
  label={lst:docente-generado-erroneo}]
// Generado por asistente computacional — CONTIENE ERRORES DE DOMINIO
public class Docente {
    private String nombre;
    private String cedula;
    private String tipoVinculacion; // ERROR-1: debería ser enum TipoDocente
    private List<String> asignaturas; // ERROR-2: String no es el tipo del dominio
    private double salario; // ERROR-3: salario no es atributo del dominio estadístico
    private Facultad facultad; // ERROR-4: cardinalidad incorrecta para cátedra

    public Docente(String nombre, String cedula, String tipoVinculacion) {
        this.nombre = nombre;
        this.cedula = cedula;
        this.tipoVinculacion = tipoVinculacion; // Sin validación del tipo
        this.asignaturas = new ArrayList<>();
    }

    public void asignarAsignatura(String codigoAsignatura) {
        // Sin verificar límite de créditos según tipo de vinculación
        this.asignaturas.add(codigoAsignatura);
    }

    public String getTipoVinculacion() {
        return tipoVinculacion;
    }

    // Getters y setters genéricos — omitidos por brevedad
}
\end{lstlisting}

El listado~\ref{lst:docente-corregido} presenta la versión corregida, producida
después de aplicar el protocolo de revisión completo. Los cuatro errores han
sido subsanados y el código es ahora coherente con el dominio del \acuerdo.

\begin{lstlisting}[style=javaStyle,
  caption={Clase \texttt{Docente} corregida mediante aplicación del protocolo
    de revisión crítica. Se emplean tipos propios del dominio, validaciones
    normativas y documentación precisa referenciando el Acuerdo 004 de 2025.},
  label={lst:docente-corregido}]
/**
 * Representa un docente vinculado a la Universidad Distrital Francisco José de Caldas.
 * La vinculación docente se clasifica según el Artículo 31 del Acuerdo 004 de 2025.
 *
 * @author Roberto Albeiro Pava-Díaz
 * @version 1.0
 */
public class Docente {
    private final String nombre;
    private final String identificacion;
    /** Tipo de vinculación: PLANTA, CATEDRA u OCASIONAL (Art. 31 Acuerdo 004/2025). */
    private final TipoDocente tipoVinculacion;
    /** Una Facultad principal para docentes de planta; puede ser null para cátedra. */
    private UnidadAcademica unidadPrincipal;
    /** Proyectos curriculares a los que está asignado este semestre. */
    private final List<ProyectoCurricular> proyectosAsignados;

    public Docente(String nombre, String identificacion, TipoDocente tipo) {
        Objects.requireNonNull(nombre, "El nombre del docente no puede ser nulo");
        Objects.requireNonNull(identificacion, "La identificación es obligatoria");
        Objects.requireNonNull(tipo, "El tipo de vinculación es obligatorio");
        this.nombre = nombre;
        this.identificacion = identificacion;
        this.tipoVinculacion = tipo;
        this.proyectosAsignados = new ArrayList<>();
    }

    /**
     * Asigna el docente a un proyecto curricular, validando la restricción
     * de carga académica máxima según el tipo de vinculación (Art. 35 Ac. 004/2025).
     *
     * @param proyecto el proyecto curricular a asignar
     * @throws IllegalStateException si se supera la carga académica permitida
     */
    public void asignarProyecto(ProyectoCurricular proyecto) {
        Objects.requireNonNull(proyecto, "El proyecto no puede ser nulo");
        int cargaActual = proyectosAsignados.stream()
            .mapToInt(ProyectoCurricular::getCreditosSemestrales)
            .sum();
        int limiteCreditos = tipoVinculacion.getLimiteCreditos();
        if (cargaActual + proyecto.getCreditosSemestrales() > limiteCreditos) {
            throw new IllegalStateException(
                String.format("Docente %s excede límite de %d créditos (%s)",
                    nombre, limiteCreditos, tipoVinculacion));
        }
        this.proyectosAsignados.add(proyecto);
    }

    public TipoDocente getTipoVinculacion() { return tipoVinculacion; }
    public String getNombre()               { return nombre; }
    public String getIdentificacion()       { return identificacion; }

    /** Devuelve la unidad académica principal; puede ser vacío para docentes de cátedra. */
    public Optional<UnidadAcademica> getUnidadPrincipal() {
        return Optional.ofNullable(unidadPrincipal);
    }
}
\end{lstlisting}

La diferencia entre los dos listados no es estilística: es de corrección semántica.
El primero compila; el segundo es correcto respecto del dominio. Esa diferencia ---entre
lo que compila y lo que es correcto--- constituye el núcleo del aprendizaje en diseño
orientado a objetos, y es, precisamente, la dimensión que el software generativo no
puede sustituir.

%% ─────────────────────────────────────────────
\section{Evaluación de la originalidad y del razonamiento en diseño orientado a objetos}
\label{sec:evaluacion-originalidad}
%% ─────────────────────────────────────────────

\noindent
La irrupción de los asistentes computacionales generativos en los entornos de
aprendizaje de programación suscita una pregunta que no es nueva pero adquiere
urgencia renovada: ¿qué se evalúa cuando se evalúa un programa? Si la pregunta
era discutible antes de 2023, hoy es inescapable. Un estudiante puede producir,
en minutos, código funcionalmente correcto para cualquier enunciado típico de
un examen de programación orientada a objetos. El código puede no ser
sintácticamente correcto, puede contener errores de dominio, puede violar SOLID ---pero, en
la mayoría de los exámenes tradicionales, esas deficiencias no se detectan.

\subsection{El reto de la evaluación auténtica}
\label{subsec:evaluacion-autentica}

La evaluación auténtica en POO ---la que mide lo que realmente importa---
no consiste en verificar que el código compila. Consiste en verificar que el
estudiante comprende por qué el código tiene la estructura que tiene, qué
alternativas de diseño se consideraron y descartaron, y bajo qué criterios
normativos o técnicos se tomó cada decisión. Esa comprensión no puede generarse
automáticamente: debe construirla el estudiante.

El marco de evaluación centrado en el razonamiento de diseño, que se propone
en la Tabla~\ref{tab:evaluacion-comp}, desplaza el foco del producto (el código)
hacia el proceso (las decisiones documentadas). Las fuentes de evidencia son
los artefactos producidos a lo largo de los doce \textit{sprints}: tarjetas de
decisión, diagramas UML comentados, reflexiones de sprint y la actividad final
del Sprint~12.

\begingroup
\small
\begin{longtable}{@{} L{2.8cm} L{3.8cm} L{4.2cm} L{2.8cm} @{} }
	% --- Título y encabezado de la primera página ---
	\caption{Comparación entre el modelo de evaluación tradicional y el modelo centrado en el razonamiento de diseño. Fuente: Elaboración del autor.} \label{tab:evaluacion-comp} \\
	\toprule
	\textbf{Aspecto} & \textbf{Evaluación tradicional} & \textbf{Evaluación centrada en razonamiento} & \textbf{Vulnerabilidad al software gen.} \\
	\midrule
	\endfirsthead
	
	% --- Encabezado para las páginas siguientes ---
	\multicolumn{4}{c}{{\bfseries \tablename\ \thetable{} -- Continuación}} \\
	\toprule
	\textbf{Aspecto} & \textbf{Evaluación tradicional} & \textbf{Evaluación centrada en razonamiento} & \textbf{Vulnerabilidad al software gen.} \\
	\midrule
	\endhead
	
	% --- Pie de tabla para páginas intermedias ---
	\midrule
	\multicolumn{4}{r}{{Continúa en la siguiente página...}} \\
	\bottomrule
	\endfoot
	
	% --- Pie de tabla final ---
	\bottomrule
	\endlastfoot
	
	% --- Datos de la tabla ---
	Qué se entrega & 
	Código fuente o examen en papel / plataforma. & 
	Código + tarjeta de decisión + diagrama UML anotado + reflexión crítica. & 
	\textbf{Alta} (solo código) / \textbf{Baja} (conjunto completo). \\
	\midrule
	
	Qué se mide & 
	Corrección sintáctica y funcionalidad básica del programa. & 
	Coherencia entre decisiones de diseño, norma estatutaria de la \ud\ y código. & 
	\textbf{Alta} / \textbf{Muy baja}. \\
	\midrule
	
	Escala temporal & 
	Evento puntual y asincrónico (examen único). & 
	Proceso longitudinal y progresivo (12 sprints). & 
	\textbf{Alta} / \textbf{Baja}. \\
	\midrule
	
	Rol del docente & 
	Calificador del producto final (corrector de errores). & 
	Evaluador del proceso, de la lógica y de la argumentación técnica. & 
	--- \\
	\midrule
	
	Referencia normativa & 
	Enunciado estático del examen o taller. & 
	\acuerdo\ y estatuto académico real de la institución. & 
	\textbf{Alta} (enunciado) / \textbf{Baja} (norma específica). \\
	\midrule
	
	Resistencia al \textit{free-riding} & 
	Baja: difícil verificar autoría en código entregado por medios digitales. & 
	Alta: la discusión oral de las tarjetas de decisión requiere conocimiento propio. & 
	--- \\
	\midrule
	
	Evidencia de aprendizaje & 
	Nota cuantitativa final. & 
	Portafolio de decisiones de diseño con evolución documentada del sistema. & 
	\textbf{Alta} / \textbf{Baja}. \\
\end{longtable}
\endgroup

\subsection{El portafolio de decisiones como instrumento de evaluación}
\label{subsec:portafolio}

El portafolio de decisiones de diseño —acumulado a lo largo de los doce
\textit{sprints} del libro— constituye el instrumento de evaluación más
resistente al software generativo, habida cuenta de que documenta no el
resultado sino el proceso deliberativo que condujo a cada elección de diseño.
Una tarjeta de decisión bien elaborada responde a cinco preguntas: ¿qué
alternativas de diseño se consideraron?, ¿cuál se eligió y por qué?, ¿qué
artículo del \acuerdo\ sustenta la elección?, ¿qué se sacrificó al elegir
esta alternativa?, y ¿qué pasaría si el requisito normativo cambiara? Ninguna
de estas preguntas admite respuesta generada automáticamente de manera
plausible sin comprensión del dominio.

La evidencia acumulada en los doce sprints incluye: diagramas UML progresivos
(Sprints~1--3, capítulos~\ref{cap:fundamentos}--\ref{cap:modelado}); código
con Javadoc específico del dominio (Sprints~4--7, capítulos~\ref{cap:java}
y~\ref{cap:rust}); suite de pruebas unitarias con invariantes normativas
(Sprints~8--9, capítulo~\ref{cap:pruebas}); reflexiones de sprint escritas;
y la reflexión crítica final del Sprint~12. Ese conjunto ---heterogéneo,
longitudinal, con trazabilidad explícita--- es lo que permite al evaluador
distinguir el razonamiento genuino del producido por delegación acrítica.

\subsection{Limitaciones del modelo propuesto}
\label{subsec:limitaciones-evaluacion}

No resulta del todo claro que el modelo de evaluación centrado en razonamiento
sea escalable en todos los contextos de la \ud. Un grupo de cuarenta estudiantes
con un docente de tiempo completo puede sostener la revisión oral de las tarjetas
de decisión; un grupo de ochenta con docentes de cátedra que dictan doce horas
semanales en tres instituciones distintas difícilmente puede hacerlo. El autor
reconoce esta limitación estructural: el modelo es pedagógicamente superior
al examen tradicional, pero impone costos de evaluación que no todas las
condiciones institucionales permiten absorber. Queda abierta la pregunta de cómo
diseñar versiones de menor costo del portafolio de decisiones sin perder su
ventaja evaluativa esencial.

%% ─────────────────────────────────────────────
\section{Marco ético para el uso responsable de software generativo}
\label{sec:marco-etico}
%% ─────────────────────────────────────────────

\noindent
La ética del software generativo en educación no es una cuestión abstracta
de filosofía moral: es un conjunto de decisiones prácticas que las instituciones,
los docentes y los estudiantes deben tomar con consecuencias reales y verificables.
El marco que se propone aquí se articula en tres dimensiones ---transparencia,
responsabilidad y equidad--- que no son independientes, sino imbricadas: la
ausencia de cualquiera de ellas menoscaba las otras dos.

\subsection{Transparencia: la declaración del uso de software generativo}
\label{subsec:transparencia}

Declarar que se empleó un asistente computacional generativo en la elaboración
de un entregable académico no es una confesión de trampa: es un acto de honestidad
intelectual. La transparencia exige que el estudiante especifique qué asistente
empleó, para qué tarea concreta, en qué momento del proceso, y qué modificaciones
realizó al resultado generado. Una declaración vaga ---«usé IA»--- no satisface
el criterio de transparencia; una declaración precisa ---«empleé Claude 3.7 para
generar el \textit{boilerplate} de los getters de la clase \texttt{UnidadAcademica},
revisé el código con el protocolo de la sección~\ref{sec:analisis-critico} y
corregí tres errores de dominio»--- sí lo satisface.

La \ud, al momento de redactar este texto, no dispone de una política institucional
explícita sobre el uso de software generativo en evaluaciones de programación.
El autor propone que esa política se incorpore al reglamento estudiantil en un
plazo razonable, estableciendo el formato de declaración de uso que se considerará
aceptable.\footnote{En la Facultad de Ingeniería de la \ud, el autor ha implementado
desde el segundo semestre de 2024 una ``declaración de autoría extendida'' que cada
entregable debe incluir, especificando las fuentes ---humanas y computacionales---
que contribuyeron al producto. La experiencia inicial sugiere que la declaración
aumenta la reflexividad del estudiante respecto de su propio proceso, aunque no
garantiza honestidad en todos los casos. El diseño de mecanismos de verificación
queda como problema abierto.}

\subsection{Responsabilidad: el estudiante es responsable del código que entrega}
\label{subsec:responsabilidad}

Con independencia del origen del código ---escrito por el estudiante, generado
por un asistente, adaptado de un tutorial, producido en colaboración con un
compañero--- el estudiante que lo entrega es responsable de su corrección y de
su coherencia con el dominio. Esta responsabilidad no es mitigable por la
procedencia del código: no existe, en términos académicos, la exculpación de
«el asistente lo generó y yo no supe que estaba mal». Saber si está bien o mal
es, precisamente, lo que el curso intenta enseñar.

Esta dimensión implica que el protocolo de revisión crítica de la
sección~\ref{sec:analisis-critico} no es opcional: es el ejercicio de la
responsabilidad que el estudiante asume al entregar código de cualquier origen.
Y si el estudiante no aplica el protocolo porque no lo conoce o porque el plazo
es insuficiente, la responsabilidad de esa condición se distribuye entre el
estudiante, el diseño del curso y la institución.

\subsection{Equidad: la brecha de acceso a software generativo}
\label{subsec:equidad}

La tercera dimensión es, probablemente, la más incómoda para una institución
pública. La diferencia entre un estudiante de la Pontificia Universidad Javeriana
o de la Universidad de los Andes ---que dispone de laptop propia con GPU dedicada,
acceso a versiones premium de múltiples asistentes generativos y conexión de
alta velocidad desde su residencia--- y un estudiante de la \ud\ que comparte
una sala de sistemas en la sede Macarena, con ancho de banda insuficiente en
horas pico y acceso limitado a versiones gratuitas de estos sistemas, no es
una diferencia de talento ni de esfuerzo: es una diferencia de infraestructura
que el mercado produce y que la institución pública tiene la obligación de
reconocer y gestionar.\footnote{El autor observó, durante el primer semestre
de 2025, que varios estudiantes de estratos 1 y 2 de la \ud\ empleaban los
kioscos de internet del barrio Santafé ---a tres cuadras de la sede Macarena---
para acceder a asistentes generativos durante horas de la tarde, puesto que el
ancho de banda de la sala de sistemas se congestionaba después de las 10 a.m.
Esta observación no constituye evidencia sistemática, pero ilustra una condición
estructural cuya invisibilización sería un error intelectual y una injusticia
práctica.}

La brecha generativa ---término empleado en la sección~\ref{subsec:penetracion-col}---
no es inevitable. Requiere respuestas institucionales concretas: provisión de
acceso a versiones educativas de asistentes generativos en las salas de sistemas,
actualización del ancho de banda, capacitación de los docentes en el uso crítico
de estos sistemas, y diseño de actividades que no presuponen acceso premium.
Ninguna de estas respuestas es técnicamente imposible; todas implican decisiones
presupuestales que la \ud\ deberá tomar con base en criterios de equidad educativa.

La Tabla~\ref{tab:marco-etico} sintetiza las tres dimensiones del marco ético.

\begingroup
\small
\begin{longtable}{@{} L{2.5cm} L{4.2cm} L{4.2cm} L{3.1cm} @{} }
	% --- Título y encabezado de la primera página ---
	\caption{Dimensiones del marco ético para el uso de software generativo en el contexto académico de la \ud. Elaboración del autor.} \label{tab:marco-etico} \\
	\toprule
	\textbf{Dimensión} & \textbf{Definición operativa} & \textbf{Ejemplo institucional} & \textbf{Mecanismo propuesto} \\
	\midrule
	\endfirsthead
	
	% --- Encabezado para las páginas siguientes ---
	\multicolumn{4}{c}{{\bfseries \tablename\ \thetable{} -- Continuación}} \\
	\toprule
	\textbf{Dimensión} & \textbf{Definición operativa} & \textbf{Ejemplo institucional} & \textbf{Mecanismo propuesto} \\
	\midrule
	\endhead
	
	% --- Pie de tabla para páginas intermedias ---
	\midrule
	\multicolumn{4}{r}{{Continúa en la siguiente página...}} \\
	\bottomrule
	\endfoot
	
	% --- Pie de tabla final ---
	\bottomrule
	\endlastfoot
	
	% --- Datos de la tabla ---
	Transparencia & 
	Declarar el uso de asistentes generativos, especificando el sistema, la tarea y las modificaciones realizadas por el autor. & 
	Inclusión en los entregables de una sección «Declaración de uso de IA» con formato estructurado por el docente. & 
	Formato de declaración estándar integrado en el reglamento de la asignatura. \\
	\midrule
	
	Responsabilidad & 
	El estudiante que entrega el código es el único responsable de su corrección, independientemente de su origen tecnológico. & 
	Si el código del Sprint~5 viola el Art.~31 del \acuerdo, el estudiante no puede excusarse en el error del asistente. & 
	Inclusión explícita del principio de autoría en el contrato pedagógico del curso. \\
	\midrule
	
	Equidad & 
	Reconocimiento institucional de la brecha de acceso y gestión activa para garantizar igualdad de oportunidades. & 
	Acceso gratuito a versiones Pro en salas de sistemas; ancho de banda garantizado para todos durante la clase. & 
	Convenio institucional con proveedores; inversión en infraestructura de red y licencias educativas. \\
\end{longtable}
\endgroup

\subsection{Recomendaciones de política para la \ud}
\label{subsec:recomendaciones-politica}

A partir del marco ético propuesto, se formulan cuatro recomendaciones concretas
de política para la Facultad de Ingeniería de la \ud, en el entendido de que
ninguna institución de educación superior puede permitirse ignorar la irrupción
de estos sistemas en los entornos de aprendizaje:

\begin{enumerate}[label=\arabic*.]
  \item Establecer una política explícita sobre el uso de software generativo
        en evaluaciones de programación, diferenciando los tipos de evaluación
        (formativa, sumativa, acreditación de competencias) y los niveles de
        autonomía permitidos.
  \item Garantizar acceso equitativo a versiones educativas de al menos un
        asistente computacional generativo de calidad comparable a los disponibles
        en universidades privadas de la ciudad, gestionado mediante convenio
        institucional.
  \item Incorporar en los planes de estudio de los programas de ingeniería
        de sistemas al menos una asignatura o componente curricular dedicado al
        análisis crítico del software generativo, sus fundamentos técnicos,
        sus limitaciones y sus implicaciones éticas.
  \item Revisar el Reglamento Estudiantil para incluir una definición precisa
        de las conductas constituyen deshonestidad académica en el contexto
        del software generativo, diferenciándolas del uso legítimo y declarado.
\end{enumerate}

%% ─────────────────────────────────────────────
\section{Comparativa \java\ vs.\ \rust: experiencia con asistentes generativos}
\label{sec:java-rust-generativo}
%% ─────────────────────────────────────────────

\noindent
Un aspecto que la literatura disponible aborda de manera insuficiente ---al menos
en lo que el autor ha podido revisar al momento de redactar este texto--- es la
diferencia cualitativa en la calidad del código generado por asistentes
computacionales para \java\ y para \rust\ en el contexto del aprendizaje
de la POO. Esa diferencia no es trivial, y tiene repercusiones directas
sobre las estrategias pedagógicas que se adoptan en las asignaturas de
la Facultad de Ingeniería de la \ud.\footnote{Las observaciones que se
documentan en esta sección son cualitativas y no constituyen evidencia
empírica controlada. Se trata de impresiones acumuladas a lo largo del
primer semestre de 2025 en dos asignaturas: Programación Orientada a Objetos
(sexto semestre) y Paradigmas de Programación (octavo semestre). El autor
advierte que la extrapolación exige cautela y que un estudio controlado
con diseño experimental riguroso podría arrojar resultados distintos.}

\subsection{La asimetría del corpus de entrenamiento}
\label{subsec:asimetria-corpus}

La razón de la diferencia cualitativa es, en principio, sencilla de enunciar:
\java\ existe desde 1995 y cuenta con más de dos décadas de código público
en repositorios, foros, tutoriales, libros y proyectos de código abierto.
El corpus de entrenamiento de los principales asistentes computacionales
generativos contiene, razonablemente, varias veces más código \java\ que
código \rust\ ---lenguaje cuya especificación estabilizada data de 2015 y cuya
adopción masiva en proyectos de código abierto se aceleró principalmente a
partir de 2020. Más corpus no garantiza mejor calidad para un dominio específico,
pero sí produce mayor familiaridad con los patrones sintácticamente correctos del lenguaje.

Pero la diferencia no es solo cuantitativa. \rust\ introduce conceptos ---el
sistema de \textit{ownership}, el \textit{borrow checker}, los tiempos de vida,
los tipos afines--- que no tienen equivalente en los lenguajes más representados
en el corpus \parencite{Klabnik2023}. El asistente generativo ha ``visto'' millones de líneas de código
con referencias nulas en \java, \texttt{C\#} y Python; ha ``visto'' considerablemente
menos código que gestiona correctamente los préstamos de referencias en \rust.
El resultado es predecible: el código \rust\ generado comete con mayor frecuencia
errores de \textit{ownership} y \textit{lifetime} que el código \java\ generado
comete errores equivalentes de gestión de referencias.

\subsection{Observaciones cualitativas del autor}
\label{subsec:observaciones-java-rust}

Cuatro patrones de error se observaron de manera recurrente en el código \rust\
generado por asistentes durante el primer semestre de 2025:

\begin{itemize}
  \item \textbf{Clonado innecesario}: el asistente resuelve conflictos de
        \textit{ownership} mediante llamadas a \texttt{.clone()} donde el
        diseño correcto emplearía préstamos (\textit{borrows}). El código
        compila, pero desperdicia memoria y viola la semántica de propiedad.

  \item \textbf{\textit{Lifetimes} explícitos omitidos}: el asistente omite
        anotaciones de tiempo de vida donde el compilador las requiere,
        produciendo errores que el estudiante novel no puede interpretar
        fácilmente.

  \item \textbf{Uso de \texttt{unwrap()} sin tratamiento de errores}:
        el asistente generativo frecuentemente sustituye el manejo
        correcto de \texttt{Result<T, E>} con llamadas a \texttt{.unwrap()},
        que producen \textit{panics} en tiempo de ejecución ante condiciones
        que el código debería manejar de manera explícita.

  \item \textbf{Estructuras con estado mutable compartido}:
        el asistente propone con frecuencia el uso de \texttt{\texttt{Rc\textless RefCell\textless T\textgreater\textgreater}}
        donde el diseño sintácticamente correcto de \rust\ emplearía patrones de propiedad
        exclusiva o mensajería mediante canales, lo que suscita código correcto
        en apariencia pero contrario al espíritu del lenguaje.
        
\end{itemize}

Para \java, los errores de dominio son más frecuentes que los errores sintacticos.
El asistente produce código que sigue las convenciones del lenguaje razonablemente
bien, pero que viola las restricciones del \acuerdo\ tal como ilustraron los
ejemplos de la sección~\ref{subsec:ejemplos-violacion}. En \rust, los errores
son tanto de dominio como sintácticamente correctos, lo que duplica la carga de revisión
crítica para el estudiante.

\subsection{Tabla comparativa de calidad del código generado}
\label{subsec:tabla-comparativa}

La Tabla~\ref{tab:java-rust-generativo} sintetiza las diferencias observadas
entre el código \java\ y \rust\ generado por asistentes computacionales.

\begingroup
\small
\begin{longtable}{@{} L{3cm} L{3.8cm} L{3.8cm} L{3.2cm} @{} }
	% --- Título y encabezado de la primera página ---
	\caption{Comparación cualitativa de la calidad del código generado para \java\ y \rust\ (Sistema Estatutario \ud). Elaboración del autor.} \label{tab:java-rust-generativo} \\
	\toprule
	\textbf{Aspecto} & \textbf{Calidad \java} & \textbf{Calidad \rust} & \textbf{Observación} \\
	\midrule
	\endfirsthead
	
	% --- Encabezado para las páginas siguientes ---
	\multicolumn{4}{c}{{\bfseries \tablename\ \thetable{} -- Continuación}} \\
	\toprule
	\textbf{Aspecto} & \textbf{Calidad \java} & \textbf{Calidad \rust} & \textbf{Observación} \\
	\midrule
	\endhead
	
	% --- Pie de tabla para páginas intermedias ---
	\midrule
	\multicolumn{4}{r}{{Continúa en la siguiente página...}} \\
	\bottomrule
	\endfoot
	
	% --- Pie de tabla final ---
	\bottomrule
	\endlastfoot
	
	% --- Datos de la tabla ---
	Idiomaticidad & 
	\textbf{Alta}: Cumple convenciones de nombres, estructura de paquetes y Javadoc estándar. & 
	\textbf{Media-baja}: Clonado innecesario (\textit{over-cloning}) y uso de \texttt{unwrap()} excesivo. & 
	Relacionado con la densidad del corpus de entrenamiento disponible para cada lenguaje. \\
	\midrule
	
	Corrección de dominio & 
	\textbf{Media}: Errores semánticos frecuentes por falta de contexto normativo local. & 
	\textbf{Media}: Los mismos errores de dominio más errores sintácticos adicionales. & 
	Ningún asistente conoce el \acuerdo\ sin que se le proporcione explícitamente en el \textit{prompt}. \\
	\midrule
	
	Manejo de errores & 
	\textbf{Correcto}: Uso convencional de bloques \texttt{try-catch}. & 
	\textbf{Incorrecto}: Uso de \texttt{unwrap()} en lugar de la propagación sintácticamente correcta con el operador \texttt{?}. & 
	El asistente tiende a minimizar la longitud del código, no a maximizar su robustez sistémica. \\
	\midrule
	
	Herencia y polimorfismo & 
	\textbf{Estable}: Bien gestionado para jerarquías simples; errores en jerarquías profundas ($>3$ niveles). & 
	\textbf{Inestable}: \textit{Traits} básicos correctos; despacho dinámico (\texttt{dyn Trait}) frecuentemente incorrecto. & 
	La complejidad del dominio (Estatuto Académico) amplifica los errores de lógica del asistente. \\
	\midrule
	
	Compilabilidad inmediata & 
	\textbf{Alta}: Estimación entre el 70\% y el 85\% de los listados generados compilan sin cambios. & 
	\textbf{Media}: Estimación entre el 40\% y el 60\%; el \textit{borrow checker} rechaza la mayoría. & 
	\textbf{Crítico}: La compilabilidad no debe confundirse con la corrección semántica o normativa. \\
\end{longtable}
\endgroup

La implicación pedagógica de esta asimetría es concreta: el docente que emplea
\rust\ en un curso de POO puede aprovechar los errores del asistente como
material de análisis, puesto que son más visibles y más frecuentes que en \java.
El error de \textit{borrow checker} que el asistente produce y que el
compilador rechaza es, paradójicamente, más pedagógicamente valioso que el
código \java\ que compila sin error pero viola el dominio: el primero es
detectable; el segundo no lo es sin el protocolo de la
sección~\ref{sec:analisis-critico}.

Y sin embargo, la asimetría plantea también un riesgo: el estudiante que trabaja
principalmente con \java\ puede desarrollar una confianza excesiva en el código
generado, precisamente porque compila y luce correcto. Esa confianza no justificada
es, en opinión del autor, el riesgo pedagógico de mayor trascendencia que el
software generativo introduce en los cursos de programación orientada a objetos.

%% ─────────────────────────────────────────────────────────────────
%% SPRINT 12
%% ─────────────────────────────────────────────────────────────────
\section*{Sprint 12: Integración, entrega final y reflexión crítica}
\label{sprint:12}
\addcontentsline{toc}{section}{Sprint 12: Integración, entrega final y reflexión crítica}

\begin{tcolorbox}[sprintBox, title={Sprint 12 --- Integración, entrega final y reflexión crítica}]
  \textbf{Duración estimada:} 2 semanas \\[4pt]
  \textbf{Objetivo general:} Integrar la totalidad de los módulos del sistema
  de información estatutario (organizacional, académico, evaluación y calidad)
  desarrollados en los Sprints~1--11; ejecutar la batería completa de pruebas
  con cobertura mínima del 80\,\%; preparar y sustentar la presentación final
  ante un panel evaluador; y producir la reflexión crítica escrita sobre el
  proceso de aprendizaje y el papel del software generativo a lo largo de los
  doce \textit{sprints}. \\[4pt]
  \textbf{Prerrequisitos:} Todos los módulos de los Sprints~1--11 compilando
  individualmente; suite de pruebas de los Sprints~8--11 pasando; diagramas UML
  finales coherentes con el código entregado.
\end{tcolorbox}

\bigskip

\subsection*{Tabla del Sprint 12}
\label{tab:sprint12-header}

\begin{longtable}{L{3.5cm} L{4.5cm} L{4.5cm}}
  \sprintcabecera
  \endhead
  %% ── Fila 1 ──
  Sistema integrado compilable y ejecutable: todos los módulos
  (organizacional, académico, evaluación) conectados mediante la
  capa de servicio, con batería completa de pruebas pasando
  ($\geq$80\,\% de cobertura según JaCoCo en \java\
  y \texttt{cargo tarpaulin} en \rust)
  & Integrar los módulos desarrollados en los Sprints~4--11 en un
    sistema coherente, verificar que las interfaces entre módulos
    respetan los contratos definidos en los diagramas UML, y
    documentar el proceso de integración en el informe técnico final
  & El sistema integrado compila sin errores, ejecuta el flujo
    completo desde la interfaz gráfica hasta la persistencia, y
    todas las pruebas de regresión pasan; la cobertura de ramas
    supera el 80\,\% \\
  \midrule
  %% ── Fila 2 ──
  Informe técnico de integración: documento de 8--12 páginas que
  describe la arquitectura final del sistema, las decisiones de
  integración tomadas, los conflictos encontrados y cómo se
  resolvieron, y las lecciones aprendidas
  & Comunicar de manera estructurada el proceso de integración,
    con énfasis en las decisiones de diseño que debieron adaptarse
    al integrar módulos desarrollados en momentos distintos del semestre
  & El informe documenta al menos tres decisiones de integración
    no triviales con argumentación técnica; referencia los diagramas
    UML de los sprints anteriores para mostrar la evolución \\
  \midrule
  %% ── Fila 3 ──
  Presentación final del proyecto (15--20 minutos) con demostración
  en vivo del sistema integrado y sustentación oral de las decisiones
  de diseño más relevantes ante un panel evaluador de dos o tres docentes
  & Demostrar comprensión profunda del sistema desarrollado, capacidad
    de argumentar las decisiones de diseño orientado a objetos y habilidad
    para responder preguntas técnicas sobre el código y el dominio
  & La presentación demuestra dominio del dominio estatutario y de los
    principios POO; las respuestas a las preguntas del panel revelan
    comprensión genuina de las decisiones de diseño y sus alternativas \\
  \midrule
  %% ── Fila 4 ──
  Portafolio de decisiones de diseño: consolidación y revisión de las
  tarjetas de decisión producidas en los doce \textit{sprints}, organizadas
  cronológicamente y con una sección de reflexión sobre la evolución del
  pensamiento de diseño del estudiante
  & Documentar la trayectoria de aprendizaje en diseño orientado a
    objetos a lo largo del semestre, identificando los momentos de
    mayor claridad conceptual y los escollos superados
  & El portafolio contiene al menos doce tarjetas de decisión no
    triviales; la sección de reflexión sobre evolución del pensamiento
    de diseño es específica y evidencia metacognición genuina \\
  \midrule
  %% ── Fila 5 ──
  Reflexión crítica escrita sobre el proceso de aprendizaje y el
  empleo de software generativo (mínimo 1.000 palabras, máximo 2.000),
  siguiendo la estructura y la rúbrica de la actividad final del Sprint~12
  & Desarrollar metacognición sobre el propio proceso de aprendizaje de
    la POO y el papel que el software generativo tuvo en él; formular
    una posición ética argumentada y una pregunta de investigación derivada
    de la experiencia propia
  & La reflexión evidencia pensamiento crítico no descriptivo; distingue
    con precisión el conocimiento propio del generado; formula una posición
    ética coherente con el marco de la sección~\ref{sec:marco-etico} y
    una pregunta de investigación específica y pertinente \\
  \midrule
  %% ── Fila 6 ──
  Versión final en \rust: adaptación del módulo organizacional y
  del módulo de evaluación al sistema de tipos de \rust, con al menos
  tres \textit{traits} propios del dominio implementados y documentados,
  y pruebas que verifiquen el comportamiento del \textit{borrow checker}
  en escenarios de propiedad compartida
  & Demostrar comprensión del modelo de propiedad de \rust\ aplicado al
    dominio estatutario, y contrastar explícitamente las decisiones de
    diseño en \rust\ con las tomadas en \java\
  & El código \rust\ es sintácticamente correcto; emplea \texttt{Result<T,E>} para
    el manejo de errores sin \texttt{unwrap()} en código de producción;
    el estudiante justifica las diferencias de diseño respecto de \java\
    con argumentación propia en las tarjetas de decisión \\
  \bottomrule
  \caption{Entregables, características y criterios de la rúbrica del Sprint~12.}
  \label{tab:sprint12}
\end{longtable}

%% ─── Rúbrica integral ────────────────────────
\subsection*{Rúbrica integral del proyecto}
\label{rubrica-integral}

\begin{longtable}{L{3cm} L{1.8cm} L{3cm} L{3cm} L{3cm}}
  \toprule
  \textbf{Dimensión} & \textbf{Peso} &
  \textbf{Insuficiente (1--2)} &
  \textbf{Aceptable (3)} &
  \textbf{Sobresaliente (4--5)} \\
  \midrule
  \endhead
  %% ── D1 ──
  Modelado UML del dominio
  & 20\,\%
  & El diagrama de clases omite entidades o relaciones del estatuto;
    las decisiones de herencia no se justifican; la multiplicidad
    de las relaciones es incorrecta o ausente
  & El diagrama es coherente con el estatuto; se justifican la mayoría
    de las decisiones de herencia y composición; la multiplicidad
    es correcta en las relaciones principales
  & El diagrama es completo y preciso; todas las decisiones se justifican
    con argumentación orientada a objetos y referencia explícita al artículo
    del \acuerdo\ aplicable; el modelo evolucionó de manera documentada
    a lo largo de los doce sprints \\
  \midrule
  %% ── D2 ──
  Implementación \java
  & 25\,\%
  & El código compila con errores o no refleja el modelo UML;
    el Javadoc es ausente o genérico y no referencia el dominio;
    los invariantes del estatuto no se verifican en el código
  & El código es funcional y coherente con el modelo; el Javadoc
    cubre los métodos principales con información específica del dominio;
    al menos las restricciones del estatuto más relevantes se validan
  & El código aplica herencia, polimorfismo y encapsulamiento de manera
    correcta; el Javadoc es preciso y referencia los artículos del
    \acuerdo\ donde corresponde; se aplican al menos dos patrones de
    diseño justificados en las tarjetas de decisión \\
  \midrule
  %% ── D3 ──
  Implementación \rust
  & 20\,\%
  & El código no compila o viola el sistema de \textit{ownership};
    emplea \texttt{unwrap()} sistemáticamente; los \textit{traits}
    implementados no corresponden al dominio
  & El código compila; los \textit{traits} propios del dominio se
    implementan correctamente; el manejo de errores emplea \texttt{Result}
    en los casos principales; el estudiante entiende los errores
    del \textit{borrow checker}
  & El código es sintácticamente correcto en \rust; no hay \texttt{unwrap()} en código
    de producción; el estudiante justifica las diferencias de diseño
    respecto de \java\ con argumentación propia y documentada \\
  \midrule
  %% ── D4 ──
  Pruebas y calidad
  & 15\,\%
  & Menos del 50\,\% de cobertura; las pruebas son triviales y no
    verifican invariantes del dominio; no hay pruebas de regresión
  & Cobertura $\geq$ 70\,\%; las pruebas verifican al menos cinco
    invariantes del dominio definidos por el \acuerdo; hay
    pruebas de casos de borde
  & Cobertura $\geq$ 80\,\%; las pruebas verifican restricciones
    normativas del estatuto con referencia explícita al artículo aplicable;
    el código refactorizado a lo largo de los sprints preserva
    el comportamiento verificado \\
  \midrule
  %% ── D5 ──
  Reflexión crítica sobre software generativo
  & 20\,\%
  & La reflexión es descriptiva o laudatoria sin análisis crítico;
    no distingue el conocimiento propio del generado; no formula
    posición ética ni pregunta de investigación
  & La reflexión identifica oportunidades y riesgos del software
    generativo con ejemplos del propio proceso; formula una posición
    ética básica; propone una pregunta de investigación pertinente
  & La reflexión evidencia pensamiento crítico profundo; contrasta
    evidencia propia con el marco ético de la
    sección~\ref{sec:marco-etico}; distingue con precisión el
    conocimiento propio del generado; formula una pregunta de
    investigación específica y derivada de la experiencia \\
  \bottomrule
  \caption{Rúbrica integral del proyecto a lo largo de los doce \textit{sprints}.}
  \label{tab:rubrica-integral}
\end{longtable}

\bigskip

\subsection*{Actividad final de reflexión sobre software generativo}
\label{act:ia-sprint12}

\begin{tcolorbox}[actividadIA,
  title={Actividad final de reflexión (Sprint~12) ---
         Software generativo y aprendizaje de la \poo}]

\textbf{Instrucción general.}
Al concluir los doce \textit{sprints} del libro, el estudiante redacta una
reflexión crítica ---mínimo 1.000 palabras, máximo 2.000--- sobre el papel
del software generativo en su proceso de aprendizaje de la programación
orientada a objetos. La reflexión no es un resumen del proyecto: es un
análisis argumentado del propio proceso, con evidencia extraída de las
tarjetas de decisión, los diagramas UML y el código producido a lo largo
del semestre. El uso de asistentes computacionales generativos para redactar
esta reflexión debe declararse explícitamente y desvirtúa el propósito del
ejercicio: lo que se evalúa es el razonamiento del estudiante sobre su propia
experiencia, no la fluidez del texto producido.

\bigskip
\textbf{Estructura sugerida (no obligatoria):}

\begin{enumerate}[label=\textbf{\arabic*.}]
  \item \textbf{Inventario de usos.}
        ¿En qué \textit{sprints} y para qué tareas específicas empleó el
        asistente computacional? Sea preciso: no «usé IA para programar»,
        sino «en el Sprint~5, empleé GitHub Copilot para generar los métodos
        \textit{getter} de las clases \texttt{Facultad} y \texttt{ProyectoCurricular},
        y Claude 3.7 para entender el error de \textit{lifetime} en la
        implementación \rust\ de la relación entre \texttt{Escuela} y
        \texttt{Docente}». El inventario debe incluir, para cada uso, si
        el resultado fue empleado directamente, modificado o descartado.

  \item \textbf{Análisis de valor.}
        ¿En qué casos el asistente generó valor genuino para el aprendizaje?
        ¿En qué casos resultó contraproducente? Proporcione al menos dos
        ejemplos de cada tipo, referenciando código o decisiones de diseño
        concretos. Un ejemplo de valor genuino puede ser la explicación de
        un mensaje de error que el estudiante no habría podido descifrar
        en el tiempo disponible. Un ejemplo de efecto contraproducente puede
        ser la aceptación de una jerarquía de herencia incorrecta que el
        asistente generó y que el estudiante no revisó con el protocolo
        de la sección~\ref{sec:analisis-critico}.

  \item \textbf{Autoevaluación del aprendizaje.}
        ¿Qué conceptos de diseño orientado a objetos comprendió genuinamente
        a lo largo de los doce \textit{sprints}? ¿Cuáles quedaron pendientes
        o se comprenden sólo superficialmente? ¿En qué medida el empleo del
        asistente computacional incidió en esa brecha? La autoevaluación debe
        ser honesta y específica: declarar que «aprendí POO» no satisface el
        criterio; declarar que «comprendí el principio de sustitución de Liskov
        en el Sprint~6 al intentar usar \texttt{DocentePlanta} donde el código
        esperaba \texttt{Docente}, y que el asistente no me ayudó a entender
        por qué eso era un error de diseño» sí lo satisface.

  \item \textbf{Posición ética.}
        ¿Cómo declararía el uso del asistente en la entrega del proyecto?
        ¿Bajo qué condiciones considera que el empleo de software generativo
        en evaluaciones de programación es éticamente aceptable? ¿Y bajo qué
        condiciones considera que constituye deshonestidad académica? La
        posición debe ser argumentada, no meramente declarativa: debe
        sustentarse en los principios del marco ético de la
        sección~\ref{sec:marco-etico} y en la experiencia propia del semestre.

  \item \textbf{Pregunta abierta de investigación.}
        Formule una pregunta de investigación sobre el empleo de software
        generativo en el aprendizaje de la programación orientada a objetos
        que usted considera relevante y no resuelta. La pregunta debe ser
        derivada de la experiencia propia, específica ---no «¿son buenos o malos
        los asistentes computacionales?»--- y pertinente para el contexto de la \ud\
        y las universidades públicas colombianas. Por ejemplo: «¿Existe
        diferencia estadísticamente significativa en la comprensión del
        principio de responsabilidad única entre estudiantes que emplearon
        software generativo para generar las clases de dominio y estudiantes
        que las escribieron manualmente, controlando por el nivel de experiencia
        previa en programación?»
\end{enumerate}

\bigskip
\textbf{Rúbrica de evaluación de la reflexión final:}

\begin{tabular}{L{3cm} L{3cm} L{3cm} L{3cm}}
  \toprule
  \textbf{Criterio} &
  \textbf{Insuficiente (1)} &
  \textbf{Aceptable (2--3)} &
  \textbf{Sobresaliente (4--5)} \\
  \midrule
  Análisis crítico
  & Meramente descriptivo; el estudiante enumera lo que hizo sin
    tomar posición argumentada sobre el valor o el riesgo del
    software generativo
  & Identifica fortalezas y limitaciones del software generativo
    con ejemplos del propio proceso; la posición es reconocible
    aunque no siempre está argumentada
  & Argumentación profunda con ejemplos del propio código y
    referencias a los conceptos del libro; la posición es
    coherente, matizada y sustentada en evidencia del semestre \\
  \midrule
  Autoevaluación del aprendizaje
  & No distingue lo que aprendió genuinamente de lo que el
    asistente computacional hizo en su lugar; no identifica
    brechas de comprensión
  & Identifica aprendizajes propios con algún nivel de especificidad
    pero sin contraste crítico sistemático con el uso del asistente;
    reconoce algunas brechas
  & Distingue con precisión el conocimiento construido del delegado;
    formula hipótesis específicas sobre el efecto del asistente
    en su aprendizaje de cada principio de POO \\
  \midrule
  Posición ética
  & Ausente, acrítica o meramente declarativa;
    no argumenta por qué su posición es la correcta
  & Formula una posición básica sobre el uso ético del software
    generativo, con alguna argumentación; reconoce que existen
    condiciones que determinan la aceptabilidad del uso
  & Formula una posición argumentada, matizada y coherente con
    el marco ético de la sección~\ref{sec:marco-etico};
    distingue tipos de uso según el nivel de formación, el tipo
    de evaluación y la disponibilidad equitativa de acceso \\
  \midrule
  Pregunta de investigación
  & Ausente, trivial o tan amplia que no es investigable;
    no se deriva de la experiencia propia
  & Pertinente y conectada a la experiencia propia pero
    demasiado amplia o con diseño de investigación no esbozado;
    podría transformarse en pregunta investigable con orientación
  & Pertinente, específica, derivada de la experiencia propia
    del semestre y con potencial de investigación empírica en
    el contexto de la \ud\ u otras universidades públicas
    colombianas; el estudiante esboza un posible diseño de estudio \\
  \bottomrule
\end{tabular}

\end{tcolorbox}

\bigskip

\subsection*{Cierre del libro: lo que permanece después del software generativo}
\label{subsec:cierre}

\noindent
Al concluir los doce \textit{sprints}, el estudiante dispone de un sistema
de información funcional, de un portafolio de decisiones de diseño, de pruebas
que verifican invariantes del \acuerdo, y ---si el proceso fue genuino--- de una
comprensión que ningún asistente computacional pudo producir por él: la intuición
de diseño que se construye al enfrentarse repetidamente al problema de modelar
un dominio complejo y al costo de las decisiones equivocadas.

Esa intuición no es transferible de un asistente a un estudiante. Se construye
iterativamente, en el error y en la revisión, en el momento en que el compilador
de \rust\ rechaza una referencia que el estudiante consideraba válida y lo obliga
a pensar en quién posee el valor. En el momento en que el revisor detecta que
la jerarquía de unidades académicas viola el artículo del estatuto y el estudiante
debe rehacer el diagrama desde la raíz. En el momento en que la prueba falla
porque el dominio tiene una restricción que el código ignoraba.

El software generativo puede acelerar algunas partes de ese proceso. Pero el
proceso mismo ---el encuentro repetido con la complejidad de un dominio real,
la negociación entre lo que el código permite y lo que la norma exige, la
búsqueda de la abstracción que hace al sistema extensible sin romper lo que ya
funciona--- ese proceso es irreductiblemente propio. Queda la duda de si la
generación siguiente de asistentes computacionales cambiará esta afirmación.
Por ahora, no resulta del todo claro que lo haga.
